\documentclass{article}
% 默认字体大小是
% Formatting
\usepackage[utf8]{inputenc}
\usepackage[margin=1in]{geometry}
\usepackage[titletoc,title]{appendix}

\usepackage{amsmath,amsfonts,amssymb,mathtools}
\usepackage{graphicx,float}
\usepackage{booktabs}
\usepackage[ruled,vlined]{algorithm2e}
\usepackage{algorithmic}
\usepackage{biblatex}
\usepackage{cleveref}
\usepackage{indentfirst}
\setlength{\parindent}{2em}

\addbibresource{references.bib}

\title{PRML Spring 2026 Homework 2: Non-Linear Supervised Learning}
\author{Yifan Ding \\
        \small Code: \url{https://github.com/douhenhuidyf/PRML-Spring-2026/tree/main/hw2} \\
        \small Email: 23373493@buaa.edu.cn}
\date{April 8, 2026}

\begin{document}
\maketitle

% \begin{abstract}

% \end{abstract}


\section{Introduction}
\subsection{Problem Statement}
Given a dataset of 1000 samples generated by a given process, categorized into two classes (C0 and C1),
 the task is to train classification models using Decision Trees, AdaBoost with Decision Trees, and Support Vector Machines (SVM) with at least three different kernel functions.
 The performance of these models will be evaluated on a test set of 500 samples, which are generated from the same distribution as the training data (250 samples for each class).
 The goal is to compare the classification performance of these algorithms and discuss why certain algorithms may perform better for this specific problem.

\subsection{Problem Analysis}
The dataset provided for this task is generated from a 3D make-moons distribution, which is known for its non-linear and intertwined structure.
This type of dataset poses a challenge for linear classifiers, as the decision boundary between the two classes (C0 and C1) is not a simple linear hyperplane.
Instead, it requires models that can capture complex relationships and non-linear patterns in the data.
Decision Trees are a natural choice for this problem, as they can model non-linear decision boundaries by recursively partitioning the feature space.
However, a single decision tree may not be sufficient to capture the complexity of the data, and it may suffer from overfitting if the tree is too deep.
AdaBoost with Decision Trees can help mitigate this issue by combining multiple weak learners (shallow trees) to create a stronger model that can better capture the underlying structure of the data.
Support Vector Machines (SVMs) with kernel functions are also well-suited for this problem, as they can implicitly map the data into a higher-dimensional feature space where a linear decision boundary can be found.
The choice of kernel function (e.g., linear, polynomial, RBF) will play a crucial role in the performance of the SVM, as it determines how well the model can capture the non-linear relationships in the data.


Overall, the non-linear nature of the dataset suggests that models capable of capturing complex decision boundaries, such as AdaBoost with Decision Trees and SVMs with appropriate kernels, are likely to perform better than a single Decision Tree or a linear SVM on this task.


\subsection{Structure of the Paper}
The paper is organized as follows: In \Cref{sec:methodology}, we describe the evaluation metrics used to assess the performance of the classification models, as well as the methodologies for Decision Trees, AdaBoost with Decision Trees, and Support Vector Machines.
In \Cref{sec:experiments}, we present the experimental setup and results for each of the three approaches.
Finally, in \Cref{sec:conclusions}, we summarize our findings and discuss the implications of our results.

\section{Methodology}
\label{sec:methodology}

\subsection{Metrics}
\label{sec:metrics}
In classification problems, common evaluation results include TP (True Positives), FP (False Positives), TN (True Negatives), and FN (False Negatives). 
Based on these results, we can calculate various performance metrics such as Accuracy, Precision, Recall, and F1 Score.
Here are the definitions of these metrics:
\begin{itemize}
        \item \textbf{Accuracy} is the ratio of correctly predicted instances to the total instances. It is calculated as:
        \[
        \text{Accuracy} = \frac{TP + TN}{TP + FP + TN + FN}
        \]
        
        \item \textbf{Precision} is the ratio of correctly predicted positive instances to the total predicted positive instances. It is calculated as:
        \[
        \text{Precision} = \frac{TP}{TP + FP}
        \]
        
        \item \textbf{Recall} (also known as Sensitivity) is the ratio of correctly predicted positive instances to the total actual positive instances. It is calculated as:
        \[
        \text{Recall} = \frac{TP}{TP + FN}
        \]
        
        \item \textbf{F1 Score} is the harmonic mean of Precision and Recall, providing a balance between the two metrics. It is calculated as:
        \[
        F1 = 2 \times \frac{\text{Precision} \times \text{Recall}}{\text{Precision} + \text{Recall}}
        \]
\end{itemize}
These metrics will be used to evaluate the performance of the classification models trained on the dataset.

\subsection{Decision Trees}
Decision Trees are a type of supervised learning algorithm that can be used for classification and regression tasks. 
They work by recursively splitting the data into multiple subsets based on the feature that provides the best separation of the classes.
In regression trees, the splitting criterion is typically based on minimizing the mean squared error (MSE) of the target variable in the resulting subsets.
While in classification trees, the splitting criterion is often based on maximizing the information gain or minimizing the Gini impurity.
Gain indicates how well a feature can separate the classes(how much probability is reduced after the split).
The recursive process continues until a stopping criterion is met, such as reaching a maximum depth or having a minimum number of samples in a leaf node.

Intuitively, Decision Trees can be visualized as splitting the feature space into multiple regions in vertical and horizontal directions, where each region corresponds to a leaf node that predicts a specific class label.
In this way, Decision Trees can capture non-linear relationships between the features and the target variable, making them suitable for classification tasks with complex decision boundaries.
However, Decision Trees can be prone to overfitting and are sensitive to small changes in the data, especially when the tree is deep and has many branches.
To mitigate these issues, techniques such as pruning, setting a maximum depth, or using ensemble methods like Random Forests and AdaBoost can be employed to improve the generalization performance of Decision Trees.

\subsection{AdaBoost with Decision Trees}
AdaBoost, short for Adaptive Boosting, is an ensemble learning method that combines multiple weak classifiers to create a strong classifier.
Specifically, AdaBoost works by iteratively training a sequence of weak classifiers, where each classifier focuses on the samples that were misclassified by the previous classifiers.
In the context of using Decision Trees as weak classifiers, AdaBoost typically uses shallow trees (often called decision stumps) that have a maximum depth of one or two.
The AdaBoost algorithm assigns weights to each training sample, which are updated after each iteration based on the performance of the current weak classifier.
The weights of the misclassified samples are increased, while the weights of the correctly classified samples are decreased, allowing the next weak classifier to focus more on the difficult samples.
The final strong classifier is a weighted combination of the weak classifiers, where the weights are determined by the performance of each weak classifier on the training data.
AdaBoost can significantly improve the performance of Decision Trees by reducing bias and variance, making it a powerful method for classification tasks, especially when the underlying decision boundary is complex and non-linear.

\subsection{Support Vector Machines}

\subsubsection{Kernel Tricks}
Kernel tricks are a powerful technique used in machine learning algorithms, particularly in support vector machines (SVMs) and kernel regression, to enable the modeling of nonlinear relationships between input features and output values without explicitly transforming the data into a higher-dimensional space.
The key idea behind kernel tricks is to use a kernel function, which computes the inner product of the data points in the original input space, to implicitly map the data into a higher-dimensional feature space where linear algorithms can be applied.
The kernel function can be defined as:
\begin{equation}
    K(x, x') = \langle \phi(x), \phi(x') \rangle,
\end{equation}
where $K(x, x')$ is the kernel function, $\phi(x)$ is the mapping function that transforms the input data into a higher-dimensional feature space, and $\langle \cdot, \cdot \rangle$ denotes the inner product in that feature space.
This allows for the capture of complex nonlinear relationships without the computational cost associated with explicitly transforming the data.

\subsubsection{Support Vector Machines}
Support Vector Machines (SVMs) leverage the concept of maximizing the margin between the decision boundary and the closest data points (support vectors) to achieve better generalization performance.
In the case of linear SVMs, the decision boundary is a hyperplane that separates the classes in the feature space. However, when the data is not linearly separable, SVMs can utilize kernel tricks to implicitly map the data into a higher-dimensional space where a linear decision boundary can be found.
Support Vector Machines with kernel functions can be modeled as follows:
\begin{equation}
\begin{aligned}
    \min_{w, b} \frac{1}{2} \|w\|^2 \\ 
    s.t. \space y_i (w^T \phi(x_i) + b) \geq 1 \space i = 1, 2, ..., N,
\end{aligned}
\end{equation}
where $w$ is the weight vector, $b$ is the bias term, $C$ is the regularization parameter, and $\xi_i$ are the slack variables.


\section{Experiments}
\label{sec:experiments}

\subsection{Experimental Setup}
We generate a 3D make-moons dataset saved in \texttt{moons\_3d.npz} with a fixed random seed.
The training set contains 500 samples and the test set contains 500 samples.
All methods are evaluated with the metrics in \Cref{sec:metrics}.

\subsection{Decision Trees}
We evaluate a single decision tree with different splitting criteria and hyper-parameters. The results are shown in \Cref{tab:dt_results}. In general, a larger maximum depth leads to a more expressive partition of the 3D feature space, improving test performance but also increasing the risk of overfitting.

\begin{table}[H]
    \centering
    \caption{Decision tree performance with different settings.}
    \label{tab:dt_results}
    \begin{tabular}{lccc|cccc|cccc}
        \toprule
        Gain & Depth & MinSamples &  & \multicolumn{4}{c}{Test} & \multicolumn{4}{c}{Train} \\
        \cmidrule(lr){5-8} \cmidrule(lr){9-12}
         &  &  &  & Acc & Prec & Rec & F1 & Acc & Prec & Rec & F1 \\
        \midrule
        Gini    & 3 & 2 &  & 0.6880 & 0.6331 & 0.8940 & 0.7413 & 0.7390 & 0.6770 & 0.9140 & 0.7779 \\
        Gini    & 5 & 2 &  & 0.8690 & 0.8727 & 0.8640 & 0.8683 & 0.9120 & 0.9274 & 0.8940 & 0.9104 \\
        Gini    & 7 & 2 &  & 0.9210 & 0.9202 & 0.9220 & 0.9211 & 0.9610 & 0.9695 & 0.9520 & 0.9606 \\
        Entropy & 5 & 2 &  & 0.7910 & 0.7782 & 0.8140 & 0.7957 & 0.8300 & 0.8149 & 0.8540 & 0.8340 \\
        \bottomrule
    \end{tabular}
\end{table}

From \Cref{tab:dt_results}, the Gini criterion performs better than entropy on this dataset. Increasing depth from 5 to 7 yields a clear improvement on both train and test sets, suggesting that the underlying decision boundary is highly non-linear and benefits from a more complex partition.

\subsection{AdaBoost with Decision Trees}
We use AdaBoost with shallow decision trees (maximum depth 2) as weak learners. The number of estimators is fixed to 50. The results are shown in \Cref{tab:adaboost_results}. Compared with a single shallow tree, AdaBoost significantly improves performance by iteratively focusing on hard-to-classify samples.

\begin{table}[H]
    \centering
    \caption{AdaBoost with depth-2 decision trees (50 estimators).}
    \label{tab:adaboost_results}
    \begin{tabular}{lcccc}
        \toprule
        Split & Acc & Prec & Rec & F1 \\
        \midrule
        Test  & 0.9270 & 0.9279 & 0.9260 & 0.9269 \\
        Train & 0.9530 & 0.9450 & 0.9620 & 0.9534 \\
        \bottomrule
    \end{tabular}
\end{table}

We observe a modest generalization gap (Train Acc 0.9530 vs Test Acc 0.9270), which is expected for an ensemble with increased capacity. Nevertheless, the test performance is strong and stable, indicating that boosting effectively reduces bias while controlling variance through shallow base learners.

\subsection{Support Vector Machines}
We evaluate SVM classifiers with three kernel functions: linear, polynomial, and RBF. All parameters are kept as the library defaults, and only the kernel type is changed. The results are summarized in \Cref{tab:svm_results}.

\begin{table}[H]
    \centering
    \caption{SVM performance with different kernels (test set).}
    \label{tab:svm_results}
    \begin{tabular}{lcccc}
        \toprule
        Kernel & Acc & Prec & Rec & F1 \\
        \midrule
        Linear & 0.6850 & 0.6755 & 0.7120 & 0.6933 \\
        Poly   & 0.8650 & 0.9551 & 0.7660 & 0.8502 \\
        RBF    & 0.9800 & 0.9819 & 0.9780 & 0.9800 \\
        \bottomrule
    \end{tabular}
\end{table}

The linear kernel performs poorly, which is consistent with the strong non-linearity of the 3D make-moons boundary. The polynomial kernel improves accuracy but shows an imbalance between precision and recall, suggesting that the chosen polynomial degree under default settings may not fully match the boundary complexity. The RBF kernel achieves the best performance, which is expected because it provides a flexible non-linear decision surface that can adapt to the curved, intertwined structure of the moons.

\section{Conclusions}
\label{sec:conclusions}

We compared three non-linear supervised learning approaches on the same 3D two-class dataset.
For decision trees, increasing depth improves performance, and the Gini criterion outperforms entropy under the tested settings.
AdaBoost with shallow trees further improves the test performance over a single tree by combining multiple weak learners.
Among SVM kernels, the RBF kernel achieves the strongest results (Acc/F1 around 0.98), indicating that this dataset benefits from smooth and highly flexible non-linear decision boundaries.


\printbibliography

% \begin{appendices}

% \section{MATLAB Functions}
% Add your important MATLAB functions here with a brief implementation explanation. This is how to make an \textbf{unordered} list:
% \begin{itemize}
%     \item \texttt{y = linspace(x1,x2,n)} returns a row vector of \texttt{n} evenly spaced points between \texttt{x1} and \texttt{x2}. 
%     \item \texttt{[X,Y] = meshgrid(x,y)} returns 2-D grid coordinates based on the coordinates contained in the vectors \texttt{x} and \texttt{y}. \text{X} is a matrix where each row is a copy of \texttt{x}, and \texttt{Y} is a matrix where each column is a copy of \texttt{y}. The grid represented by the coordinates \texttt{X} and \texttt{Y} has \texttt{length(y)} rows and \texttt{length(x)} columns.  
% \end{itemize}

% \end{appendices}

\end{document}
